\documentclass[11pt,a4paper]{article}

% ---------- Packages ----------
\usepackage[margin=1in]{geometry}
\usepackage{amsmath, amssymb, amsthm, mathtools}
\usepackage{microtype}
\usepackage{algorithm}
\usepackage{algpseudocode}
\usepackage[colorlinks=true,
            linkcolor=blue!50!black,
            urlcolor=blue!50!black,
            citecolor=blue!50!black]{hyperref}

% ---------- Theorem environments ----------
\theoremstyle{plain}
\newtheorem{theorem}{Theorem}[section]
\newtheorem{proposition}[theorem]{Proposition}
\newtheorem{lemma}[theorem]{Lemma}
\newtheorem{corollary}[theorem]{Corollary}

\theoremstyle{definition}
\newtheorem{definition}[theorem]{Definition}

\theoremstyle{remark}
\newtheorem{remark}[theorem]{Remark}
\newtheorem{example}[theorem]{Example}

% ---------- Math shortcuts ----------
\newcommand{\R}{\mathbb{R}}
\newcommand{\N}{\mathbb{N}}
\newcommand{\E}{\mathbb{E}}
\newcommand{\Var}{\operatorname{Var}}
\newcommand{\Cov}{\operatorname{Cov}}
\newcommand{\Corr}{\operatorname{Corr}}
\newcommand{\F}{\mathcal{F}}
\newcommand{\Prob}{\mathbb{P}}
\newcommand{\Q}{\mathbb{Q}}
\newcommand{\dd}{\mathop{}\!\mathrm{d}}
\newcommand{\eqdef}{\overset{\mathrm{def}}{=}}
\newcommand{\indic}{\mathbf{1}}

% ---------- Document metadata ----------
\title{Variance Reduction in Monte Carlo:\\
       Theoretical Foundations}
\author{Quant Finance Portfolio --- Phase 2, Block 2.0}
\date{\today}

\begin{document}
\maketitle

\begin{abstract}
We develop the two foundational variance reduction techniques used
throughout the rest of Phase~2: antithetic variates and control
variates. For each we derive the conditions under which it is
guaranteed to reduce variance, the resulting variance formula, and
the practical algorithm that implements it. The treatment is honest
about the cost-benefit trade-off: a method that halves the variance
but doubles the per-path cost is no improvement, and we make the
\emph{variance times time} (VxT) criterion explicit at the outset to
guard against this confusion. We also derive the empirical
control-variate coefficient and analyse the small bias it introduces,
and we close with a brief catalogue of techniques we do not cover in
the present block (stratified sampling, importance sampling,
multilevel MC) together with the reasons.
\end{abstract}

\tableofcontents

\section{Introduction}
\label{sec:intro}

Every Monte Carlo estimator we have built so far has had the same
asymptotic structure: $\hat I_M = M^{-1} \sum_{i=1}^M Y_i$ for some
i.i.d.\ payoff $Y$, with variance $\Var(\hat I_M) = \sigma^2_Y / M$
where $\sigma^2_Y = \Var(Y)$. The half-width of a fixed-confidence
interval scales as $\sigma_Y / \sqrt{M}$.

For a typical at-the-money European call (parameters $S_0 = K = 100$,
$r = 0.05$, $\sigma = 0.20$, $T = 1$) the payoff variance is
$\sigma^2_Y \approx 217$. A half-width of $0.01$ at $95\%$ confidence
requires $M \approx 4 \cdot 1.96^2 \cdot 217 / 0.01^2 \approx 8.3
\cdot 10^6$ paths. Each ten-fold reduction in the half-width costs a
hundred-fold increase in the sample size.

Variance reduction techniques attack this rate by constructing
\emph{alternative} estimators of the same quantity $I = \E[Y]$, with
the same sample size $M$, but with variance smaller than $\sigma^2_Y
/ M$. The reduction comes from exploiting structure of the problem
(monotonicity, correlation with a tractable quantity) that the basic
estimator ignores.

The general framework is: replace $Y$ with another random variable
$\tilde Y$ such that $\E[\tilde Y] = \E[Y]$ but $\Var(\tilde Y) <
\Var(Y)$, and run the Monte Carlo with $\tilde Y$ in place of $Y$.
The estimator $\hat I_M^{\tilde Y} = M^{-1} \sum_i \tilde Y_i$ has
variance $\Var(\tilde Y) / M$. If $\Var(\tilde Y) = \sigma^2_Y / k$,
the variance is reduced by a factor of $k$, which is equivalent to
multiplying the effective sample size by $k$.

The catch, almost always, is that producing one sample of $\tilde Y$
costs more than producing one sample of $Y$. The honest way of
comparing is the \emph{variance times time} criterion of
Section~\ref{sec:vxt}.

\section{The variance-times-time (VxT) criterion}
\label{sec:vxt}

Suppose two estimators of the same quantity $I$ are available:
\begin{itemize}
\item Estimator $A$: each sample has variance $\sigma_A^2$ and takes
time $\tau_A$ to compute.
\item Estimator $B$: each sample has variance $\sigma_B^2$ and takes
time $\tau_B$ to compute.
\end{itemize}
Given a total computational budget $T_{\text{total}}$, estimator $A$
can produce $M_A = T_{\text{total}} / \tau_A$ samples, leading to
estimator variance $\sigma_A^2 / M_A = \sigma_A^2 \tau_A /
T_{\text{total}}$. Similarly, $B$ produces estimator variance
$\sigma_B^2 \tau_B / T_{\text{total}}$.

The estimator with smaller \emph{product} $\sigma^2 \tau$ wins, and
the relative VxT,
\begin{equation}
\frac{\sigma_A^2 \tau_A}{\sigma_B^2 \tau_B},
\end{equation}
quantifies how many times faster $B$ converges to the truth at fixed
budget. A variance reduction technique that halves the variance but
doubles the per-sample cost is a wash: VxT unchanged, no actual
improvement.

This criterion will reappear in every validation script of Phase~2
Block~2. Variance ratios alone are misleading; only the VxT ratio
captures the practical comparison.

\begin{remark}
The VxT criterion assumes the per-sample cost is the same for every
sample, which is true for the simple estimators considered here. For
adaptive estimators (where the cost varies with the realisation), a
more careful analysis is needed; see Glasserman~\cite{Glasserman},
Section~4.1.3.
\end{remark}

\section{Antithetic variates}
\label{sec:av}

\subsection{Setup and theorem}

Suppose $Y$ depends on a vector of standard normals $\mathbf{Z}$:
$Y = f(\mathbf{Z})$ with $\mathbf{Z} \sim \mathcal{N}(0, I)$. By
symmetry of the normal distribution, $-\mathbf{Z}$ is also
$\mathcal{N}(0, I)$, so $f(-\mathbf{Z})$ has the same distribution
(and the same expectation) as $f(\mathbf{Z})$.

The antithetic variates estimator pairs the two:
\begin{equation}
\label{eq:av-Yi}
Y_i^{\mathrm{AV}} \;\eqdef\; \tfrac{1}{2}\bigl(\,f(\mathbf{Z}_i) + f(-\mathbf{Z}_i)\,\bigr),
\qquad i = 1, \ldots, M,
\end{equation}
with $\mathbf{Z}_i$ i.i.d.\ standard normals across $i$. The estimator
of $I = \E[f(\mathbf{Z})]$ is
\begin{equation}
\hat I_M^{\mathrm{AV}} \;\eqdef\; \frac{1}{M} \sum_{i=1}^M Y_i^{\mathrm{AV}}.
\end{equation}

\begin{theorem}[Variance of the antithetic estimator]
\label{thm:av}
The estimator $\hat I_M^{\mathrm{AV}}$ is unbiased for $I$, and its
variance equals
\begin{equation}
\label{eq:av-variance}
\Var\bigl(\hat I_M^{\mathrm{AV}}\bigr)
\;=\; \frac{1}{2M} \bigl(\,\Var(f(\mathbf{Z})) + \Cov(f(\mathbf{Z}),\, f(-\mathbf{Z}))\,\bigr).
\end{equation}
In particular, $\Var(\hat I_M^{\mathrm{AV}}) < \Var(\hat I_M^{\mathrm{IID}})$
if and only if $\Cov(f(\mathbf{Z}), f(-\mathbf{Z})) <
\Var(f(\mathbf{Z}))$, equivalently $\Corr(f(\mathbf{Z}), f(-\mathbf{Z}))
< 1$, which is the generic case.
\end{theorem}

\begin{proof}
Unbiasedness is immediate: $\E[Y_i^{\mathrm{AV}}] = \tfrac{1}{2}(\E[f(\mathbf{Z})]
+ \E[f(-\mathbf{Z})]) = \E[f(\mathbf{Z})] = I$. For the variance,
$\Var(Y_i^{\mathrm{AV}}) = \tfrac{1}{4}(\Var(f(\mathbf{Z})) +
\Var(f(-\mathbf{Z})) + 2 \Cov(f(\mathbf{Z}), f(-\mathbf{Z})))$, and
$\Var(f(-\mathbf{Z})) = \Var(f(\mathbf{Z}))$ by distributional
equality, so $\Var(Y_i^{\mathrm{AV}}) = \tfrac{1}{2}(\Var(f(\mathbf{Z}))
+ \Cov(f(\mathbf{Z}), f(-\mathbf{Z})))$. The estimator's variance is
this divided by $M$.
\end{proof}

The strict comparison with the IID estimator is more informative when
made at \emph{equal computational cost}, since each $Y_i^{\mathrm{AV}}$
requires two evaluations of $f$ rather than one. The honest
comparison is between $\hat I_M^{\mathrm{AV}}$ and $\hat I_{2M}^{\mathrm{IID}}$,
both using $2M$ payoff evaluations:
\begin{equation}
\Var\bigl(\hat I_{2M}^{\mathrm{IID}}\bigr)
\;=\; \frac{\Var(f(\mathbf{Z}))}{2M}.
\end{equation}
The AV estimator wins if and only if $\Cov(f(\mathbf{Z}), f(-\mathbf{Z}))
< 0$. This brings us to the key sufficient condition.

\subsection{The monotonicity condition}

\begin{theorem}[Sufficient condition for AV to reduce variance]
\label{thm:av-monotone}
If $f : \R^d \to \R$ is monotone in each coordinate of $\mathbf{Z}$
(in the same direction, i.e.\ either monotone increasing in all or
monotone decreasing in all), then
\begin{equation}
\Cov(f(\mathbf{Z}),\, f(-\mathbf{Z})) \leq 0,
\end{equation}
with strict inequality unless $f$ is essentially constant.
\end{theorem}

This is a special case of a general result on monotone functions of
random vectors and antithetic transformations; see
Glasserman~\cite[Section~4.2]{Glasserman}. We do not reproduce the
proof; the intuition is that if $f$ is monotone, large values of
$f(\mathbf{Z})$ correspond to favourable directions of $\mathbf{Z}$,
and reflecting $\mathbf{Z}$ to $-\mathbf{Z}$ flips these to
unfavourable directions, producing small values of $f(-\mathbf{Z})$.
The pair $(f(\mathbf{Z}), f(-\mathbf{Z}))$ is therefore negatively
correlated.

For the European call under GBM, $S_T = S_0 \exp((r - \sigma^2/2) T
+ \sigma \sqrt{T} Z)$ is strictly increasing in $Z$, and the call
payoff $(S_T - K)^+$ is monotone increasing in $S_T$, hence in $Z$.
So the monotonicity condition holds and AV is guaranteed to reduce
variance for the European call. By how much depends on how strong the
negative correlation is, which we measure empirically in Block~2.1.

\section{Control variates}
\label{sec:cv}

\subsection{Setup and theorem}

Suppose, in addition to the target $Y$, we have access to a second
random variable $X$ defined on the same paths, with the property that
$\E[X]$ is \emph{known in closed form}. The variable $X$ is called a
\emph{control variate}.

For any constant $c \in \R$, the centred quantity $X - \E[X]$ has
mean zero, so the random variable
\begin{equation}
\label{eq:cv-Yi}
\tilde Y(c) \;\eqdef\; Y - c\,(X - \E[X])
\end{equation}
satisfies $\E[\tilde Y(c)] = \E[Y] - c \cdot 0 = \E[Y] = I$. The
estimator built from $\tilde Y(c)$ is unbiased for $I$ regardless of
$c$. The freedom to choose $c$ lets us reduce variance.

\begin{theorem}[Optimal control coefficient]
\label{thm:cv-optimal}
The variance of $\tilde Y(c)$ is
\begin{equation}
\Var(\tilde Y(c)) \;=\; \Var(Y) - 2c\,\Cov(Y, X) + c^2\,\Var(X),
\end{equation}
minimised at
\begin{equation}
\label{eq:cv-c-star}
c^* \;\eqdef\; \frac{\Cov(Y, X)}{\Var(X)},
\end{equation}
with minimum value
\begin{equation}
\label{eq:cv-min-var}
\Var(\tilde Y(c^*)) \;=\; \Var(Y)\,(1 - \rho^2),
\qquad \rho \;\eqdef\; \Corr(Y, X).
\end{equation}
\end{theorem}

\begin{proof}
The variance formula is a straightforward expansion of $\tilde
Y(c)$. Setting the derivative with respect to $c$ to zero yields
$c^* = \Cov(Y, X) / \Var(X)$. Substituting back,
$\Var(\tilde Y(c^*)) = \Var(Y) - \Cov(Y, X)^2 / \Var(X) = \Var(Y)(1
- \rho^2)$ since $\rho^2 = \Cov(Y, X)^2 / (\Var(Y) \Var(X))$.
\end{proof}

The variance reduction factor $1 / (1 - \rho^2)$ depends only on the
correlation $\rho$ between $Y$ and the control $X$. A control with
$|\rho| = 0.9$ reduces variance by a factor $1/(1-0.81) \approx
5.3$; with $|\rho| = 0.99$, by $\approx 50$. Finding a good control
is a creative problem: it must (a) have known expectation in closed
form and (b) be highly correlated with $Y$.

\subsection{Connection with OLS regression}

The optimal coefficient $c^*$ has a direct interpretation in terms
of ordinary least squares. Consider regressing $Y$ on $X$ with an
intercept:
\begin{equation}
Y \;=\; \alpha + \beta X + \varepsilon,
\end{equation}
with $\E[\varepsilon] = 0$ and $\Cov(\varepsilon, X) = 0$. The OLS
slope estimator (in the population) is exactly
\begin{equation}
\beta^{\text{OLS}} \;=\; \frac{\Cov(Y, X)}{\Var(X)} \;=\; c^*.
\end{equation}
The control variate adjustment $\tilde Y(c^*) = Y - c^*(X - \E[X])$
is the \emph{residual} of the OLS regression, shifted by the
intercept $\alpha = \E[Y] - c^* \E[X]$ to recover the original mean.
The variance of the residual is exactly $\Var(Y)(1 - R^2)$ with
$R^2 = \rho^2$, recovering~\eqref{eq:cv-min-var}.

This is more than a calculation curiosity: it tells us that
\emph{control variate variance reduction is just OLS regression
applied to Monte Carlo}. Multiple control variates correspond to
multiple regression, and the same intuition (more controls help if
they bring genuine new information beyond the existing controls,
i.e.\ they are not collinear with them) applies. We do not pursue
multiple controls in this block, but the framework extends
mechanically.

\subsection{The empirical \texorpdfstring{$c^*$}{c-star} and its bias}
\label{subsec:empirical-c}

In practice $\Cov(Y, X)$ and $\Var(X)$ are unknown. We estimate them
from the same $M$ paths used for the Monte Carlo:
\begin{equation}
\hat c \;\eqdef\;
\frac{\hat\Cov(Y, X)}{\hat\Var(X)}
\;=\; \frac{\sum_{i=1}^M (Y_i - \bar Y)(X_i - \bar X)}
            {\sum_{i=1}^M (X_i - \bar X)^2},
\end{equation}
which is exactly the OLS slope estimate from the sample. The
estimator becomes
\begin{equation}
\hat I_M^{\mathrm{CV}} \;\eqdef\; \frac{1}{M} \sum_{i=1}^M
\bigl(Y_i - \hat c\,(X_i - \E[X])\bigr).
\end{equation}

This estimator is no longer strictly unbiased for $I$. Because
$\hat c$ is itself a function of $\{(Y_i, X_i)\}_{i=1}^M$, its
correlation with the $Y_i - \hat c\,(X_i - \E[X])$ terms induces a
small bias.

\begin{proposition}[Bias of the empirical CV estimator]
\label{prop:cv-bias}
Under standard moment conditions, the bias of $\hat I_M^{\mathrm{CV}}$
is $O(1/M)$, while its variance is $\Var(Y)(1-\rho^2)/M + O(1/M^2)$.
\end{proposition}

The proof is a perturbation analysis: $\hat c = c^* + O_p(M^{-1/2})$,
and the bias contribution at this order vanishes by orthogonality.
See Glasserman~\cite[Section~4.1.3]{Glasserman} for details.

The practical consequence: the bias decays at the same rate as the
square root of the standard error decays, so for any $M$ at which
the standard error is acceptably small, the bias is also negligibly
small (in fact, asymptotically negligible compared to the standard
error itself, since bias is $O(1/M)$ while standard error is
$O(1/\sqrt M)$). For $M \geq 10^3$ in our typical setting, the bias
is below $10^{-4}$, well within rounding tolerance for option
pricing.

The alternative is a \emph{pilot run} of $M_0$ paths to estimate $c^*$
followed by a production run of $M$ paths with $\hat c$ frozen, which
gives an exactly unbiased estimator. We do not adopt this in our
implementation: the empirical $\hat c$ approach is simpler, the bias
is negligible at our sample sizes, and the OLS interpretation is
more transparent.

\subsection{The two controls we use}

For the European call under GBM, two natural controls present
themselves:

\paragraph{Control 1 -- the underlying:} $X_1 \eqdef e^{-rT} S_T$, the
discounted terminal price. Its expectation is $\E[X_1] = e^{-rT} \E[S_T]
= e^{-rT} \cdot S_0 e^{rT} = S_0$, by the risk-neutral martingale
property. This control is correlated with the call payoff $(S_T -
K)^+$ but only modestly: the correlation breaks down where the
payoff is zero (in the OTM region), so the linear relationship is
weak. We expect $\rho \approx 0.5$--$0.7$ at the ATM point, giving
a variance reduction factor of $\approx 1.4$--$2.0$.

\paragraph{Control 2 -- the asset-or-nothing:} $X_2 \eqdef e^{-rT}\,
S_T \indic_{\{S_T > K\}}$, the discounted asset-or-nothing payoff.
Its expectation has a closed Black-Scholes form
\begin{equation}
\E[X_2] \;=\; S_0\, \Phi(d_1),
\end{equation}
with $d_1 = (\log(S_0/K) + (r + \sigma^2/2)T) / (\sigma \sqrt T)$ as
in the standard Black-Scholes formula. (This is the value of the
asset-or-nothing call, which is a building block of the
Black-Scholes derivation.) Its key property: \emph{it activates on
exactly the same paths as the call}, so its correlation with $(S_T -
K)^+$ is much higher --- typically $|\rho| > 0.95$. We expect a
variance reduction factor of $\geq 10$.

The disparity between the two controls illustrates a general
principle: \emph{a good control variate matches the ``shape'' of the
target}. The call payoff is zero on $\{S_T < K\}$ and linear in
$S_T$ on $\{S_T \geq K\}$. The discounted asset matches the linear
part but is nonzero everywhere. The asset-or-nothing matches both
the zero region and the linear region.

\section{When to use which technique}
\label{sec:when}

The two techniques apply in different circumstances and are not
mutually exclusive. A pragmatic decision tree:

\begin{itemize}
\item \emph{Is the payoff a monotone function of the underlying
randomness?} If yes, antithetic variates are essentially free
insurance: they cost one extra payoff evaluation per pair (so per-VxT
cost factor 2) but always reduce variance for monotone payoffs.

\item \emph{Is there a tractable control variate with high
correlation?} If yes, control variates can deliver order-of-magnitude
reductions, far beyond what AV achieves. The challenge is finding
the control: it requires knowing the closed-form expectation of
something correlated with the target.

\item \emph{Both?} The two can be combined: first apply antithetic
variates to construct $Y^{\mathrm{AV}}$, then use a control variate
on $Y^{\mathrm{AV}}$. The combination is implemented in production
quant systems but adds complexity; we do not pursue it in this
portfolio.
\end{itemize}

For the European call under GBM, which is our running example, the
asset-or-nothing control is so effective that AV is a marginal
addition on top. We implement and validate both separately to
illustrate the principles, then leave the combination as an
exercise.

\section{What we do not cover}
\label{sec:not-covered}

For completeness, three other variance reduction techniques exist
and are well-developed in the literature:

\paragraph{Stratified sampling.} Partition the support of an input
random variable into strata, sample a fixed (proportional) number
from each stratum, and combine the strata estimators. Reduces the
``between-strata'' component of variance but not the ``within-strata''
component. Particularly powerful in low dimensions and for inputs
where good strata can be designed (such as sampling from a single
normal: stratifying its quantile space delivers near-deterministic
convergence). See Glasserman~\cite[Chapter~4]{Glasserman}.

\paragraph{Importance sampling.} Change the probability measure to
oversample from regions where the payoff is large, then reweight to
recover an unbiased estimator. Mathematically a Radon--Nikodym
change of measure. Most effective when the payoff has a localised
``important'' region (deep out-of-the-money options, rare events).
See Glasserman~\cite[Chapter~4]{Glasserman}, Asmussen \&
Glynn~\cite{AsmussenGlynn}.

\paragraph{Multilevel Monte Carlo (MLMC).} Combine simulations at
multiple discretisation grids, exploiting that fine-grid simulations
are expensive but coarse-grid simulations are noisy. Reduces the
total cost to achieve a target accuracy for SDE-based MC. See
Giles~\cite{Giles2008}.

We do not cover these in the portfolio for reasons of scope: the
asset-or-nothing control variate alone delivers more variance
reduction than this entire course needs, and the techniques above
become valuable in regimes (deep OTM, very long-dated, exotic
path-dependent) that we do not target here.

\section{Bibliographic notes}
\label{sec:biblio}

The textbook reference for variance reduction in financial Monte
Carlo is Glasserman~\cite{Glasserman}, Chapter~4. The OLS connection
of control variates is standard, well-treated in any book on
simulation theory, e.g.\ Asmussen \& Glynn~\cite{AsmussenGlynn},
Chapter~V. The asset-or-nothing as a Black-Scholes building block is
in Hull~\cite{Hull}, Chapter~26. The general framework of
variance-times-time as the right comparison criterion is
folklore in the simulation community and stated explicitly in
Glasserman~\cite[Section~1.1.3]{Glasserman}.

\begin{thebibliography}{99}

\bibitem{AsmussenGlynn}
S.\ Asmussen and P.\ W.\ Glynn,
\emph{Stochastic Simulation: Algorithms and Analysis},
Springer, 2007.

\bibitem{Giles2008}
M.\ B.\ Giles,
\emph{Multilevel Monte Carlo path simulation},
Operations Research, 56(3):607--617, 2008.

\bibitem{Glasserman}
P.\ Glasserman,
\emph{Monte Carlo Methods in Financial Engineering},
Springer, 2004.

\bibitem{Hull}
J.\ C.\ Hull,
\emph{Options, Futures, and Other Derivatives},
10th edition, Pearson, 2018.

\end{thebibliography}

\end{document}
